%% ============================================================
%%  SOLUCIÓN — Ejercicio 4.1: Mini-presentación con Beamer
%%  Clase 2 — Después de diapositiva "Estructura básica de Beamer"
%%  Compilar: F5 (pdflatex, una o dos veces)
%% ============================================================

\documentclass[aspectratio=169]{beamer}
\usetheme{metropolis}
\usepackage[utf8]{inputenc}
\usepackage[spanish]{babel}
\usepackage{amsmath, booktabs}

% --- PERSONALIZAR COLORES ---
\definecolor{miazul}{HTML}{1B3A6B}
\definecolor{mioroado}{HTML}{D4A017}

\setbeamercolor{frametitle}{bg=miazul, fg=white}
\setbeamercolor{progress bar}{fg=mioroado}
\setbeamercolor{title separator}{fg=mioroado}

\title{Retornos a la educación en Paraguay}
\subtitle{Evidencia de la EPH 2023}
\author{Tu Nombre}
\institute{Tu Institución}
\date{\today}

\begin{document}

\maketitle

\begin{frame}{Agenda}
  \tableofcontents
\end{frame}

% ============================================================
\section{Motivación}
% ============================================================

\begin{frame}{¿Por qué estudiar los retornos a la educación?}
  \begin{itemize}
    \item \textbf{Pregunta central en economía:} ¿Cuánto ganan las personas 
          si invierten en educación?
    \item \textbf{Importancia de política:} Entender retornos ayuda a diseñar 
          políticas educativas efectivas.
    \item \textbf{Disparidades:} En Paraguay poco se conoce sobre estos retornos, 
          especialmente por género y región.
  \end{itemize}

  \vspace{1em}
  \alert{Pregunta de investigación:} ¿Cuál es el efecto causal de un año 
  adicional de educación sobre el salario mensual de los trabajadores?
\end{frame}

% ============================================================
\section{El modelo}
% ============================================================

\begin{frame}{Ecuación de Mincer}
  Nuestro modelo base es la ecuación de Mincer:
  \[
    \ln(w_i) = \alpha + \beta_1 \, \text{educ}_i + \beta_2 \, \text{exp}_i 
               + \beta_3 \, \text{exp}_i^2 + \varepsilon_i
  \]

  \begin{columns}[T]
    \begin{column}{0.48\textwidth}
      \textbf{Interpretación}
      \begin{itemize}
        \item $\beta_1$: retorno porcentual por año de educación
        \item $\beta_2, \beta_3$: ciclo de vida de ingresos (forma de U invertida)
      \end{itemize}
    \end{column}
    \begin{column}{0.48\textwidth}
      \textbf{Supuesto clave}
      \[
        \mathbb{E}[\varepsilon_i | \text{educ}_i, \text{exp}_i] = 0
      \]
      Exogeneidad: educación y experiencia no correlacionadas con el error
    \end{column}
  \end{columns}
\end{frame}

% ============================================================
\section{Datos}
% ============================================================

\begin{frame}{Fuente: EPH 2023}
  \begin{columns}[T]
    \begin{column}{0.55\textwidth}
      \textbf{Encuesta Permanente de Hogares (EPH)}
      \begin{itemize}
        \item Realizada por el INE
        \item $N = 1.200$ individuos ocupados
        \item Rango de edad: 18--65 años
        \item Representatividad: nacional
      \end{itemize}
    \end{column}
    \begin{column}{0.42\textwidth}
      \textbf{Estadísticas clave}
      \begin{center}
      \begin{tabular}{lr}
        \toprule
        Variable      & Media \\
        \midrule
        Educación     & 12.3 años \\
        Experiencia   & 8.7 años  \\
        Mujer (\%)    & 48\%      \\
        \bottomrule
      \end{tabular}
      \end{center}
    \end{column}
  \end{columns}
\end{frame}

% ============================================================
\section{Resultados}
% ============================================================

\begin{frame}{Resultados principales}
  \begin{center}
  \begin{tabular}{lccc}
    \toprule
                  & (1)      & (2)      & (3)      \\
    \midrule
    Educación     & 0.082*** & 0.071*** & 0.069*** \\
                  & (0.008)  & (0.009)  & (0.010)  \\
    Experiencia   &          & 0.034**  & 0.031**  \\
                  &          & (0.012)  & (0.013)  \\
    Mujer (=1)    &          &          & $-$0.183*** \\
                  &          &          & (0.042)  \\
    \midrule
    $R^2$ aj.     & 0.241    & 0.318    & 0.374    \\
    \bottomrule
  \end{tabular}
  \end{center}
  \vspace{0.5em}
  \footnotesize Errores estándar robustos en paréntesis. 
  *** $p<0.01$, ** $p<0.05$, * $p<0.1$.
\end{frame}

\begin{frame}{Interpretación}
  \begin{block}{Hallazgo 1: Retorno a la educación}
    Un año adicional de educación aumenta el salario en \textbf{6.9\%},
    controlando por experiencia y género.
  \end{block}

  \begin{alertblock}{Hallazgo 2: Brecha de género}
    Las mujeres ganan \textbf{18.3\% menos} que los hombres con 
    igual educación y experiencia.
  \end{alertblock}

  \begin{block}{Hallazgo 3: Ciclo de vida}
    El coeficiente negativo en experiencia$^2$ indica retornos 
    decrecientes a la experiencia, consistente con teoría.
  \end{block}
\end{frame}

% ============================================================
\section{Conclusiones}
% ============================================================

\begin{frame}{Conclusiones y agenda de investigación}
  \begin{columns}[T]
    \begin{column}{0.5\textwidth}
      \textbf{Hallazgos principales}
      \begin{itemize}
        \item Retorno a la educación: $\sim$7\% por año
        \item Brecha de género: $-$18.3\%
        \item Retornos robustos a especificaciones
      \end{itemize}
    \end{column}
    \begin{column}{0.47\textwidth}
      \textbf{Próximos pasos}
      \begin{itemize}
        \item Explorar heterogeneidad por región
        \item Variables instrumentales para endogeneidad
        \item Investigar fuentes de brecha de género
        \item Datos de panel para efectos fijos
      \end{itemize}
    \end{column}
  \end{columns}

  \vspace{1.5em}
  \begin{center}
    \large \textbf{¡Gracias por su atención!} \\[0.3em]
    \normalsize Preguntas y comentarios bienvenidos
  \end{center}
\end{frame}

\end{document}
